\section{Backend Technologies}
This section outlines the technology stack selected for the backend development, detailing the rationale behind each choice.

\subsection{ASP.NET Core}
ASP.NET Core \cite{aspnetcore} is a high-performance, open-source framework engineered for building modern, cloud-based, and cross-platform applications. Serving as the evolution of the traditional ASP.NET framework, it introduces a modular architecture that emphasizes speed, scalability, and developer productivity.

Our selection of ASP.NET Core was driven by the following factors:
\begin{itemize}
\item \textbf{Skill Acquisition}: We identified this framework as an ideal opportunity to expand our technical repertoire. Although not part of our formal curriculum, its widespread industry adoption presented a compelling challenge for professional growth.
\item \textbf{Career Alignment}: Given its significant market share within enterprise environments, mastering the Microsoft ecosystem provides a strategic advantage for our future career prospects.
\item \textbf{High-Performance Runtime}: The framework's lightweight, modular design ensures optimal execution speed, which is a critical requirement for a responsive digital wallet application.
\item \textbf{Cross-Platform Versatility}: Its ability to operate seamlessly across Windows, Linux, and macOS offers us the deployment flexibility necessary for diverse infrastructure requirements.
\end{itemize}

\subsection{Entity Framework Core}
Entity Framework Core (EF Core) \cite{entityframeworkcore} is the lightweight and extensible object-relational mapper (ORM) for .NET. It simplifies data access by allowing developers to interact with the underlying database using strongly-typed C\# objects. By abstracting raw SQL queries, EF Core promotes cleaner, maintainable code while providing a robust interface for complex database operations.



/*
\subsection{Swagger and Swashbuckle}

The project implementation benefited significantly from Swashbuckle\cite{swashbuckle}, a library that integrates the Swagger/OpenAPI\cite{swagger} specification into ASP.NET Core.

This library eliminated the need to manually write OpenAPI documentation instead of investing time in crafting detailed YAML specifications, Swashbuckle automatically generated interactive API documentation directly from our code. We also activated the option of xml comments to enhance the quality and clarity of the documentation.

By integrating Swashbuckle into our ASP.NET Core application, we ensured that documentation remained synchronized with the actual implementation. This consistent documentation proved particularly valuable during parallel development.
*/

\subsection{Mapster}
Mapster is a high-performance mapping library for .NET applications\cite{mapster}. It provides an efficient way to transform data between different object types without writing repetitive mapping code.

The decision to use Mapster over other mapping libraries like AutoMapper was driven by several key advantages:
\begin{itemize}
  \item Higher Performance: Mapster is 4x faster than other mapping libraries.
  \item Minimal configuration
\end{itemize}

In our application, Mapster plays a crucial role in transforming data between our domain models and DTOs in our API controllers without the need to write repetitive mapping code.

To illustrate the simplification provided by Mapster, consider the following comparison. Manually mapping properties from a list of domain objects to a list of DTOs would typically involve iterating through the list and assigning each property individually, as shown in Listing \ref{lst:mapster}.

\begin{lstlisting}[caption={Mapping without Mapster}, label={lst:mapster}]
  var domainObjects = await _service.GetAllItems(limit, offset);
  var dtoList = new List<ItemDto>();
  foreach (var domainObject in domainObjects)
  {
      dtoList.Add(new ItemDto
      {
          Id = domainObject.Id,
          Name = domainObject.Name,
          // ... map other properties manually
      });
  }
  return Ok(new { items = dtoList });
  \end{lstlisting}
  

This manual approach becomes verbose and error-prone as the number of properties and types increases. With Mapster, the same transformation can be achieved concisely:

\begin{lstlisting}[caption={Mapping with Mapster}, label={lst:mapster-adapt}]
  var domainObjects = await _service.GetAllItems(limit, offset);
  return Ok(new { items = domainObjects.Adapt<List<ItemDto>>() });
  \end{lstlisting}
  

Mapster's `Adapt` method handles the mapping automatically based on convention or configured rules, significantly reducing the amount of code required and improving readability.
/*
\subsection{Quartz.NET}
Quartz.Net is an open-source job scheduling library for .NET applications. It allows tasks to run automatically at specified times or intervals \cite{quartznet}.

We used Quartz.NET in our application to address the following critical issues:
\begin{itemize}
  \item The automatic renewal of user-defined spending limits.
  \item To automatically process recurring transactions that users have configured in the system.
\end{itemize}

Our system allows users to establish \textbf{monthly} or \textbf{weekly} spending limits on their wallets, requiring a mechanism to reset these limits at regular intervals without manual intervention.

We implemented a scheduled job to run at 00h01 every Sunday for the renewal of weekly limits because the weekly limits are based on calendar weeks, and a monthly job to run at 00h01 on the first day of every month for the renewal of monthly limits. See more details on the schedulers implementation in Section \ref{sec:schedulers}.

This was achieved using CRON scheduling expressions present on \cite{cronSchedule} and \cite{cronSchedule2}, and a function present on the database that will verify if the limit has expired and if so, it will be renewed.

The decision to implement scheduling logic in the API layer using Quartz.NET, rather than PostgreSQL's native scheduling capabilities, was driven by practical considerations. Centralizing scheduling logic within the application enhances maintainability, as all scheduling rules remain in one place, making them easier to modify and debug.

Additionally, this approach keeps business logic within the application layer, ensuring consistency and reducing dependencies on database-specific features. This makes the system more adaptable to future changes and easier to maintain.
*/

\subsection{MailKit}
MailKit is an open-source cross-platform .NET library for SMTP, POP3, and IMAP. It provides a simple way to handle email functionality within .NET applications \cite{mailkit}.

In our Digital Wallet application, MailKit is important for enhancing security and user verification when a user tries to register on the app it will send a confirmation email to the user's email address. The implementation leverages MailKit's SmtpClient to connect to our email account, authenticate, and send the confirmation code to the user's email address.

This automated email verification system is essential for maintaining the security of our application, preventing fraudulent accounts, and ensuring that all users have valid contact information in the system.
/* ??
\subsection{XUnit}

XUnit is a widely used testing framework for .NET applications\cite{xunit} providing simplicity and extensibility. Unlike older testing frameworks, XUnit uses attributes instead of base classes, enabling more flexible test organization and execution.

Our testing strategy implemented a multi-layered approach, using the Xunit Fixture pattern\cite{xunitFixture}:
\begin{itemize}
  \item \textbf{Unit Tests}: We developed targeted tests for the Service Layer, isolating business logic from external dependencies through mocking. This approach allowed us to verify business rules and service behaviors independently of infrastructure concerns.
  
  \item \textbf{Integration Tests}: For the Repository Layer, we created integration tests that verified our data access code against a real database environment. These tests ensured that our Entity Framework queries and database operations functioned correctly with the actual PostgreSQL schema.
\end{itemize}

For integration testing, we used Testcontainers\cite{testcontainers}, a library that programmatically manages Docker containers. This allowed us to:
\begin{itemize}
  \item Automatically create isolated PostgreSQL containers for each test class
  \item Pre-populate databases with controlled test data through our seeding methods
  \item Execute tests within a real database environment without affecting development or production databases
  \item Ensure consistent test results regardless of the development environment
\end{itemize}




